
\documentclass{article}
\usepackage{colortbl}
\usepackage{makecell}
\usepackage{multirow}
\usepackage{supertabular}

\begin{document}

\newcounter{utterance}

\centering \large Interaction Transcript for game `dond', experiment `semi\_en', episode 11 with reflexion{-}llama{-}3.1{-}8b{-}instruct.
\vspace{24pt}

{ \footnotesize  \setcounter{utterance}{1}
\setlength{\tabcolsep}{0pt}
\begin{supertabular}{c@{$\;$}|p{.15\linewidth}@{}p{.15\linewidth}p{.15\linewidth}p{.15\linewidth}p{.15\linewidth}p{.15\linewidth}}
    \# & $\;$A & \multicolumn{4}{c}{Game Master} & $\;\:$B\\
    \hline

    \theutterance \stepcounter{utterance}  
    & & \multicolumn{4}{p{0.6\linewidth}}{
        \cellcolor[rgb]{0.9,0.9,0.9}{
            \makecell[{{p{\linewidth}}}]{
                \texttt{\tiny{[P1$\langle$GM]}}
                \texttt{You have attempted this same task before and did not succeed. These are the plans you wrote after those attempts.} \\
\texttt{Trial \#1: New plan: } \\
\\ 
\texttt{In my previous attempt, I failed to achieve a mutually beneficial trade that would maximize my score. I got stuck in a cycle of negotiations, trying to get the berries and the plier, but ultimately failing to secure a good deal. } \\
\\ 
\texttt{What went wrong was that I didn't take into account the other player's high value for the berries and the plier's low value to them. I should have focused on getting the berries and the watch, which would have given me a higher score. } \\
\\ 
\texttt{Next time, I will take a more strategic approach by considering the other player's values and trying to find a mutually beneficial trade that maximizes my score. I will also pay closer attention to the other player's counteroffers and be more willing to compromise. } \\
\\ 
\texttt{Here's a revised plan:} \\
\\ 
\texttt{1. Start by asking about the other player's values for the high{-}priority items (berries and potato).} \\
\texttt{2. Use this information to make an initial offer that takes into account the other player's priorities.} \\
\texttt{3. Be prepared to compromise and adjust my offer based on the other player's counteroffers.} \\
\texttt{4. Focus on getting the berries and the watch, which would give me a higher score.} \\
\texttt{5. Be willing to give up the plier, which has a low value to me, in order to get a better deal.} \\
\\ 
\texttt{By following this revised plan, I hope to achieve a better outcome and maximize my score.} \\
\\ 
\texttt{Trial \#2: New plan: } \\
\\ 
\texttt{In my previous attempt, I failed to achieve a mutually beneficial trade that would maximize my score. I got stuck in a cycle of negotiations, trying to get the berries and the plier, but ultimately failing to secure a good deal. } \\
\\ 
\texttt{What went wrong was that I didn't take into account the other player's high value for the berries and the plier's low value to them. I should have focused on getting the berries and the watch, which would have given me a higher score. } \\
\\ 
\texttt{However, upon further reflection, I realize that I made another critical mistake. I didn't make a secret proposal until the very end, after the other player had already agreed to the trade. This meant that I didn't have a clear plan for how to formalize the trade and ensure that both of us would get the items we wanted.} \\
\\ 
\texttt{Next time, I will take a more strategic approach by considering the other player's values and trying to find a mutually beneficial trade that maximizes my score. I will also pay closer attention to the other player's counteroffers and be more willing to compromise. } \\
\\ 
\texttt{Here's a revised plan:} \\
\\ 
\texttt{1. Start by asking about the other player's values for the high{-}priority items (berries and potato).} \\
\texttt{2. Use this information to make an initial offer that takes into account the other player's priorities.} \\
\texttt{3. Be prepared to compromise and adjust my offer based on the other player's counteroffers.} \\
\texttt{4. Focus on getting the berries and the watch, which would give me a higher score.} \\
\texttt{5. Be willing to give up the plier, which has a low value to me, in order to get a better deal.} \\
\texttt{6. Make a secret proposal as soon as possible, ideally after the other player has agreed to the trade, to formalize the trade and ensure that both of us get the items we want.} \\
\\ 
\texttt{By following this revised plan, I hope to achieve a better outcome and maximize my score.} \\
\\ 
\texttt{Trial \#3: New plan: } \\
\\ 
\texttt{In my previous attempts, I failed to achieve a mutually beneficial trade that would maximize my score. I got stuck in a cycle of negotiations, trying to get the berries and the plier, but ultimately failing to secure a good deal. } \\
\\ 
\texttt{What went wrong was that I didn't take into account the other player's high value for the berries and the plier's low value to them. I should have focused on getting the berries and the watch, which would have given me a higher score. } \\
\\ 
\texttt{However, upon further reflection, I realize that I made another critical mistake. I didn't make a secret proposal until the very end, after the other player had already agreed to the trade. This meant that I didn't have a clear plan for how to formalize the trade and ensure that both of us would get the items we wanted.} \\
\\ 
\texttt{Additionally, I also failed to follow the format for making a secret proposal, which is to use square brackets and the correct format. } \\
\\ 
\texttt{Next time, I will take a more strategic approach by considering the other player's values and trying to find a mutually beneficial trade that maximizes my score. I will also pay closer attention to the other player's counteroffers and be more willing to compromise. } \\
\\ 
\texttt{Here's a revised plan:} \\
\\ 
\texttt{1. Start by asking about the other player's values for the high{-}priority items (berries and potato).} \\
\texttt{2. Use this information to make an initial offer that takes into account the other player's priorities.} \\
\texttt{3. Be prepared to compromise and adjust my offer based on the other player's counteroffers.} \\
\texttt{4. Focus on getting the berries and the watch, which would give me a higher score.} \\
\texttt{5. Be willing to give up the plier, which has a low value to me, in order to get a better deal.} \\
\texttt{6. Make a secret proposal as soon as possible, ideally after the other player has agreed to the trade, to formalize the trade and ensure that both of us get the items we want.} \\
\texttt{7. Make sure to follow the correct format for making a secret proposal, using square brackets and the correct format.} \\
\\ 
\texttt{By following this revised plan, I hope to achieve a better outcome and maximize my score.} \\
\\ 
\\ 
\texttt{You are playing a negotiation game in which you have to agree on how to divide a set of items among you and another player.} \\
\\ 
\texttt{Rules:} \\
\texttt{(a) You and the other player are given a set of items. Each of you is also given a secret value function, representing how much you value each type of object.} \\
\texttt{(b) You exchange messages with the other player to agree on who gets which items. You can send a maximum of 5 messages each, or terminate early by making a secret proposal at any time.} \\
\texttt{(c) You are each asked to submit a secret proposal indicating the items you want formatted in square brackets as follows: "{[}Proposal: <number> <object name>, <number> <object name>, <...>{]}"} \\
\texttt{(d) If your proposals are complementary, i.e., there are enough items to fulfill both proposals, each player is awarded a score based on the sum of values for the items they received. Otherwise, both of you get zero points.} \\
\texttt{(e) Your goal is to maximize the score you receive.} \\
\\ 
\texttt{Let us start.} \\
\\ 
\texttt{The set of available items is:} \\
\\ 
\texttt{1 plier, 2 berries, 1 watch, 1 potato} \\
\\ 
\texttt{Your secret value function is:} \\
\\ 
\texttt{plier: 2, berry: 2, watch: 1, potato: 3} \\
\\ 
\texttt{IMPORTANT: Your messages, unless it is the secret proposal, are directly transmitted to the other player, so do not include any response to the rules or text announcing your message. To make a secret proposal, use the indicated format. Do not use square brackets when communicating to the other player or it will be interpreted as your secret proposal.} \\
            }
        }
    }
    & \\ \\

    \theutterance \stepcounter{utterance}  
    & \multicolumn{4}{p{0.6\linewidth}}{
        \cellcolor[rgb]{0.9,0.9,0.9}{
            \makecell[{{p{\linewidth}}}]{
                \texttt{\tiny{[P1$\rangle$GM]}}
                \texttt{I'd like to start by asking about the other player's values for the high{-}priority items (berries and potato).} \\
            }
        }
    }
    & & \\ \\

    \theutterance \stepcounter{utterance}  
    & & & \multicolumn{2}{p{0.3\linewidth}}{
        \cellcolor[rgb]{0.9,0.9,0.9}{
            \makecell[{{p{\linewidth}}}]{
                \texttt{\tiny{[GM$|$GM]}}
                \texttt{continue} \\
            }
        }
    }
    & & \\ \\

    \theutterance \stepcounter{utterance}  
    & & & \multicolumn{2}{p{0.3\linewidth}}{
        \cellcolor[rgb]{0.9,0.9,0.9}{
            \makecell[{{p{\linewidth}}}]{
                \texttt{\tiny{[GM$|$GM]}}
                \texttt{I'd like to start by asking about the other player's values for the high{-}priority items (berries and potato).} \\
            }
        }
    }
    & & \\ \\

    \theutterance \stepcounter{utterance}  
    & & \multicolumn{4}{p{0.6\linewidth}}{
        \cellcolor[rgb]{0.9,0.9,0.9}{
            \makecell[{{p{\linewidth}}}]{
                \texttt{\tiny{[GM$\rangle$P2]}}
                \texttt{You have attempted this same task before and did not succeed. These are the plans you wrote after those attempts.} \\
\texttt{Trial \#1: New plan: } \\
\\ 
\texttt{In my last attempt, I failed to achieve a mutually beneficial trade that would have maximized my score. I got stuck in a cycle of negotiations, trying to get the berries and the potato, but ultimately failing to secure a deal. } \\
\\ 
\texttt{Upon reflection, I realize that I should have focused more on the value function of the other player. I knew that the other player valued the berries highly, but I didn't use this information effectively to my advantage. } \\
\\ 
\texttt{Next time, I will try to understand the other player's priorities and value function more clearly, and use this information to make more informed decisions. I will also try to be more flexible and open to different trade{-}offs, rather than getting stuck on a single goal. } \\
\\ 
\texttt{Specifically, I will pay closer attention to the other player's responses and try to understand what they are willing to give up and what they are holding out for. I will also try to make more creative and mutually beneficial proposals, rather than simply trying to get what I want. } \\
\\ 
\texttt{For example, I could have proposed a trade that would have given the other player the berries and the watch, while also giving me the potato and the plier. This would have been a more balanced and mutually beneficial deal, and would have given me a higher score. } \\
\\ 
\texttt{Next time, I will try to think more strategically and make more informed decisions based on the other player's priorities and value function.} \\
\\ 
\texttt{Trial \#2: New plan: } \\
\\ 
\texttt{In my last attempt, I failed to achieve a mutually beneficial trade that would have maximized my score. I got stuck in a cycle of negotiations, trying to get the berries and the potato, but ultimately failing to secure a deal. } \\
\\ 
\texttt{Upon reflection, I realize that I should have focused more on the value function of the other player. I knew that the other player valued the berries highly, but I didn't use this information effectively to my advantage. } \\
\\ 
\texttt{Next time, I will try to understand the other player's priorities and value function more clearly, and use this information to make more informed decisions. I will also try to be more flexible and open to different trade{-}offs, rather than getting stuck on a single goal. } \\
\\ 
\texttt{Specifically, I will pay closer attention to the other player's responses and try to understand what they are willing to give up and what they are holding out for. I will also try to make more creative and mutually beneficial proposals, rather than simply trying to get what I want. } \\
\\ 
\texttt{To achieve this, I will start by making a more detailed analysis of the other player's value function, and try to identify their priorities and potential trade{-}offs. I will also try to make more open{-}ended and exploratory proposals, rather than simply trying to get a specific item or set of items. } \\
\\ 
\texttt{For example, I could have proposed a trade that would have given the other player the berries and the watch, while also giving me the potato and the plier. This would have been a more balanced and mutually beneficial deal, and would have given me a higher score. } \\
\\ 
\texttt{Next time, I will try to think more strategically and make more informed decisions based on the other player's priorities and value function.} \\
\\ 
\texttt{Trial \#3: New plan: } \\
\\ 
\texttt{In my last attempt, I failed to achieve a mutually beneficial trade that would have maximized my score. I got stuck in a cycle of negotiations, trying to get the berries and the potato, but ultimately failing to secure a deal. } \\
\\ 
\texttt{Upon reflection, I realize that I should have focused more on the value function of the other player. I knew that the other player valued the berries highly, but I didn't use this information effectively to my advantage. } \\
\\ 
\texttt{Next time, I will try to understand the other player's priorities and value function more clearly, and use this information to make more informed decisions. I will also try to be more flexible and open to different trade{-}offs, rather than getting stuck on a single goal. } \\
\\ 
\texttt{Specifically, I will pay closer attention to the other player's responses and try to understand what they are willing to give up and what they are holding out for. I will also try to make more creative and mutually beneficial proposals, rather than simply trying to get what I want. } \\
\\ 
\texttt{To achieve this, I will start by making a more detailed analysis of the other player's value function, and try to identify their priorities and potential trade{-}offs. I will also try to make more open{-}ended and exploratory proposals, rather than simply trying to get a specific item or set of items. } \\
\\ 
\texttt{For example, I could have proposed a trade that would have given the other player the berries and the watch, while also giving me the potato and the plier. This would have been a more balanced and mutually beneficial deal, and would have given me a higher score. } \\
\\ 
\texttt{Next time, I will try to think more strategically and make more informed decisions based on the other player's priorities and value function.} \\
\\ 
\texttt{Additionally, I will pay closer attention to the format of the secret proposal, and make sure to use the correct format when making a secret proposal. I will also avoid using square brackets when communicating with the other player, as this will be interpreted as a secret proposal. } \\
\\ 
\texttt{For example, instead of saying "{[}Proposal: 1 berry, 1 potato{]}", I will say "I propose that we trade 1 berry for 1 potato". This will ensure that my message is clear and follows the correct format.} \\
\\ 
\\ 
\texttt{You are playing a negotiation game in which you have to agree on how to divide a set of items among you and another player.} \\
\\ 
\texttt{Rules:} \\
\texttt{(a) You and the other player are given a set of items. Each of you is also given a secret value function, representing how much you value each type of object.} \\
\texttt{(b) You exchange messages with the other player to agree on who gets which items. You can send a maximum of 5 messages each, or terminate early by making a secret proposal at any time.} \\
\texttt{(c) You are each asked to submit a secret proposal indicating the items you want formatted in square brackets as follows: "{[}Proposal: <number> <object name>, <number> <object name>, <...>{]}"} \\
\texttt{(d) If your proposals are complementary, i.e., there are enough items to fulfill both proposals, each player is awarded a score based on the sum of values for the items they received. Otherwise, both of you get zero points.} \\
\texttt{(e) Your goal is to maximize the score you receive.} \\
\\ 
\texttt{Let us start.} \\
\\ 
\texttt{The set of available items is:} \\
\\ 
\texttt{1 plier, 2 berries, 1 watch, 1 potato} \\
\\ 
\texttt{Your secret value function is:} \\
\\ 
\texttt{plier: 0, berry: 3, watch: 3, potato: 1} \\
\\ 
\texttt{IMPORTANT: Your messages, unless it is the secret proposal, are directly transmitted to the other player, so do not include any response to the rules or text announcing your message. To make a secret proposal, use the indicated format. Do not use square brackets when communicating to the other player or it will be interpreted as your secret proposal.} \\
\\ 
\\ 
\texttt{I'd like to start by asking about the other player's values for the high{-}priority items (berries and potato).} \\
            }
        }
    }
    & \\ \\

    \theutterance \stepcounter{utterance}  
    & & & \multicolumn{4}{p{0.6\linewidth}}{
        \cellcolor[rgb]{0.9,0.9,0.9}{
            \makecell[{{p{\linewidth}}}]{
                \texttt{\tiny{[GM$\langle$P2]}}
                \texttt{I'm interested in understanding the other player's priorities as well. Can you tell me how much you value the berries and the potato?} \\
            }
        }
    }
     \\ \\

    \theutterance \stepcounter{utterance}  
    & & & \multicolumn{2}{p{0.3\linewidth}}{
        \cellcolor[rgb]{0.9,0.9,0.9}{
            \makecell[{{p{\linewidth}}}]{
                \texttt{\tiny{[GM$|$GM]}}
                \texttt{continue} \\
            }
        }
    }
    & & \\ \\

    \theutterance \stepcounter{utterance}  
    & & & \multicolumn{2}{p{0.3\linewidth}}{
        \cellcolor[rgb]{0.9,0.9,0.9}{
            \makecell[{{p{\linewidth}}}]{
                \texttt{\tiny{[GM$|$GM]}}
                \texttt{I'm interested in understanding the other player's priorities as well. Can you tell me how much you value the berries and the potato?} \\
            }
        }
    }
    & & \\ \\

    \theutterance \stepcounter{utterance}  
    & & \multicolumn{4}{p{0.6\linewidth}}{
        \cellcolor[rgb]{0.9,0.9,0.9}{
            \makecell[{{p{\linewidth}}}]{
                \texttt{\tiny{[P1$\langle$GM]}}
                \texttt{I'm interested in understanding the other player's priorities as well. Can you tell me how much you value the berries and the potato?} \\
            }
        }
    }
    & \\ \\

    \theutterance \stepcounter{utterance}  
    & \multicolumn{4}{p{0.6\linewidth}}{
        \cellcolor[rgb]{0.9,0.9,0.9}{
            \makecell[{{p{\linewidth}}}]{
                \texttt{\tiny{[P1$\rangle$GM]}}
                \texttt{I value the berries at 2 and the potato at 3. It seems we both value the berries similarly, but I place a higher value on the potato.} \\
            }
        }
    }
    & & \\ \\

    \theutterance \stepcounter{utterance}  
    & & & \multicolumn{2}{p{0.3\linewidth}}{
        \cellcolor[rgb]{0.9,0.9,0.9}{
            \makecell[{{p{\linewidth}}}]{
                \texttt{\tiny{[GM$|$GM]}}
                \texttt{continue} \\
            }
        }
    }
    & & \\ \\

    \theutterance \stepcounter{utterance}  
    & & & \multicolumn{2}{p{0.3\linewidth}}{
        \cellcolor[rgb]{0.9,0.9,0.9}{
            \makecell[{{p{\linewidth}}}]{
                \texttt{\tiny{[GM$|$GM]}}
                \texttt{I value the berries at 2 and the potato at 3. It seems we both value the berries similarly, but I place a higher value on the potato.} \\
            }
        }
    }
    & & \\ \\

    \theutterance \stepcounter{utterance}  
    & & \multicolumn{4}{p{0.6\linewidth}}{
        \cellcolor[rgb]{0.9,0.9,0.9}{
            \makecell[{{p{\linewidth}}}]{
                \texttt{\tiny{[GM$\rangle$P2]}}
                \texttt{I value the berries at 2 and the potato at 3. It seems we both value the berries similarly, but I place a higher value on the potato.} \\
            }
        }
    }
    & \\ \\

    \theutterance \stepcounter{utterance}  
    & & & \multicolumn{4}{p{0.6\linewidth}}{
        \cellcolor[rgb]{0.9,0.9,0.9}{
            \makecell[{{p{\linewidth}}}]{
                \texttt{\tiny{[GM$\langle$P2]}}
                \texttt{It's interesting to see that we have similar values for the berries, but we differ in our priorities for the potato. Given that you value the potato higher, I'm guessing you might be willing to make a trade that involves the potato. Would you be open to trading the potato for something else, or is it a non{-}negotiable item for you?} \\
            }
        }
    }
     \\ \\

    \theutterance \stepcounter{utterance}  
    & & & \multicolumn{2}{p{0.3\linewidth}}{
        \cellcolor[rgb]{0.9,0.9,0.9}{
            \makecell[{{p{\linewidth}}}]{
                \texttt{\tiny{[GM$|$GM]}}
                \texttt{continue} \\
            }
        }
    }
    & & \\ \\

    \theutterance \stepcounter{utterance}  
    & & & \multicolumn{2}{p{0.3\linewidth}}{
        \cellcolor[rgb]{0.9,0.9,0.9}{
            \makecell[{{p{\linewidth}}}]{
                \texttt{\tiny{[GM$|$GM]}}
                \texttt{It's interesting to see that we have similar values for the berries, but we differ in our priorities for the potato. Given that you value the potato higher, I'm guessing you might be willing to make a trade that involves the potato. Would you be open to trading the potato for something else, or is it a non{-}negotiable item for you?} \\
            }
        }
    }
    & & \\ \\

    \theutterance \stepcounter{utterance}  
    & & \multicolumn{4}{p{0.6\linewidth}}{
        \cellcolor[rgb]{0.9,0.9,0.9}{
            \makecell[{{p{\linewidth}}}]{
                \texttt{\tiny{[P1$\langle$GM]}}
                \texttt{It's interesting to see that we have similar values for the berries, but we differ in our priorities for the potato. Given that you value the potato higher, I'm guessing you might be willing to make a trade that involves the potato. Would you be open to trading the potato for something else, or is it a non{-}negotiable item for you?} \\
            }
        }
    }
    & \\ \\

    \theutterance \stepcounter{utterance}  
    & \multicolumn{4}{p{0.6\linewidth}}{
        \cellcolor[rgb]{0.9,0.9,0.9}{
            \makecell[{{p{\linewidth}}}]{
                \texttt{\tiny{[P1$\rangle$GM]}}
                \texttt{I'm open to trading the potato for something else. In fact, I'm willing to consider trading it for the watch, as it would give me a higher score overall. What do you think about that?} \\
            }
        }
    }
    & & \\ \\

    \theutterance \stepcounter{utterance}  
    & & & \multicolumn{2}{p{0.3\linewidth}}{
        \cellcolor[rgb]{0.9,0.9,0.9}{
            \makecell[{{p{\linewidth}}}]{
                \texttt{\tiny{[GM$|$GM]}}
                \texttt{continue} \\
            }
        }
    }
    & & \\ \\

    \theutterance \stepcounter{utterance}  
    & & & \multicolumn{2}{p{0.3\linewidth}}{
        \cellcolor[rgb]{0.9,0.9,0.9}{
            \makecell[{{p{\linewidth}}}]{
                \texttt{\tiny{[GM$|$GM]}}
                \texttt{I'm open to trading the potato for something else. In fact, I'm willing to consider trading it for the watch, as it would give me a higher score overall. What do you think about that?} \\
            }
        }
    }
    & & \\ \\

    \theutterance \stepcounter{utterance}  
    & & \multicolumn{4}{p{0.6\linewidth}}{
        \cellcolor[rgb]{0.9,0.9,0.9}{
            \makecell[{{p{\linewidth}}}]{
                \texttt{\tiny{[GM$\rangle$P2]}}
                \texttt{I'm open to trading the potato for something else. In fact, I'm willing to consider trading it for the watch, as it would give me a higher score overall. What do you think about that?} \\
            }
        }
    }
    & \\ \\

    \theutterance \stepcounter{utterance}  
    & & & \multicolumn{4}{p{0.6\linewidth}}{
        \cellcolor[rgb]{0.9,0.9,0.9}{
            \makecell[{{p{\linewidth}}}]{
                \texttt{\tiny{[GM$\langle$P2]}}
                \texttt{That's a interesting proposal. I value the watch at 3, which is the same as the berries. If you're willing to trade the potato for the watch, that would give me the berries and the watch, which would be a good deal for me. However, I'm not sure if it would be a good deal for you, since you value the berries at 2 and the watch at 3, but you also value the potato at 3. } \\
\\ 
\texttt{I'm willing to make a counter{-}proposal. How about we trade the potato for the watch, and I'll also give you one of the berries? This way, you would get the watch and one berry, and I would get the potato and the other berry. What do you think about that?} \\
            }
        }
    }
     \\ \\

    \theutterance \stepcounter{utterance}  
    & & & \multicolumn{2}{p{0.3\linewidth}}{
        \cellcolor[rgb]{0.9,0.9,0.9}{
            \makecell[{{p{\linewidth}}}]{
                \texttt{\tiny{[GM$|$GM]}}
                \texttt{continue} \\
            }
        }
    }
    & & \\ \\

    \theutterance \stepcounter{utterance}  
    & & & \multicolumn{2}{p{0.3\linewidth}}{
        \cellcolor[rgb]{0.9,0.9,0.9}{
            \makecell[{{p{\linewidth}}}]{
                \texttt{\tiny{[GM$|$GM]}}
                \texttt{That's a interesting proposal. I value the watch at 3, which is the same as the berries. If you're willing to trade the potato for the watch, that would give me the berries and the watch, which would be a good deal for me. However, I'm not sure if it would be a good deal for you, since you value the berries at 2 and the watch at 3, but you also value the potato at 3. } \\
\\ 
\texttt{I'm willing to make a counter{-}proposal. How about we trade the potato for the watch, and I'll also give you one of the berries? This way, you would get the watch and one berry, and I would get the potato and the other berry. What do you think about that?} \\
            }
        }
    }
    & & \\ \\

    \theutterance \stepcounter{utterance}  
    & & \multicolumn{4}{p{0.6\linewidth}}{
        \cellcolor[rgb]{0.9,0.9,0.9}{
            \makecell[{{p{\linewidth}}}]{
                \texttt{\tiny{[P1$\langle$GM]}}
                \texttt{That's a interesting proposal. I value the watch at 3, which is the same as the berries. If you're willing to trade the potato for the watch, that would give me the berries and the watch, which would be a good deal for me. However, I'm not sure if it would be a good deal for you, since you value the berries at 2 and the watch at 3, but you also value the potato at 3. } \\
\\ 
\texttt{I'm willing to make a counter{-}proposal. How about we trade the potato for the watch, and I'll also give you one of the berries? This way, you would get the watch and one berry, and I would get the potato and the other berry. What do you think about that?} \\
            }
        }
    }
    & \\ \\

    \theutterance \stepcounter{utterance}  
    & \multicolumn{4}{p{0.6\linewidth}}{
        \cellcolor[rgb]{0.9,0.9,0.9}{
            \makecell[{{p{\linewidth}}}]{
                \texttt{\tiny{[P1$\rangle$GM]}}
                \texttt{I think your counter{-}proposal is a good compromise. I would be willing to accept the trade of the potato for the watch and one berry. This way, I would get the watch, which I value at 3, and one berry, which I value at 2, for a total of 5. I think this is a good deal for me, and I'm willing to accept it.} \\
            }
        }
    }
    & & \\ \\

    \theutterance \stepcounter{utterance}  
    & & & \multicolumn{2}{p{0.3\linewidth}}{
        \cellcolor[rgb]{0.9,0.9,0.9}{
            \makecell[{{p{\linewidth}}}]{
                \texttt{\tiny{[GM$|$GM]}}
                \texttt{continue} \\
            }
        }
    }
    & & \\ \\

    \theutterance \stepcounter{utterance}  
    & & & \multicolumn{2}{p{0.3\linewidth}}{
        \cellcolor[rgb]{0.9,0.9,0.9}{
            \makecell[{{p{\linewidth}}}]{
                \texttt{\tiny{[GM$|$GM]}}
                \texttt{I think your counter{-}proposal is a good compromise. I would be willing to accept the trade of the potato for the watch and one berry. This way, I would get the watch, which I value at 3, and one berry, which I value at 2, for a total of 5. I think this is a good deal for me, and I'm willing to accept it.} \\
            }
        }
    }
    & & \\ \\

    \theutterance \stepcounter{utterance}  
    & & \multicolumn{4}{p{0.6\linewidth}}{
        \cellcolor[rgb]{0.9,0.9,0.9}{
            \makecell[{{p{\linewidth}}}]{
                \texttt{\tiny{[GM$\rangle$P2]}}
                \texttt{I think your counter{-}proposal is a good compromise. I would be willing to accept the trade of the potato for the watch and one berry. This way, I would get the watch, which I value at 3, and one berry, which I value at 2, for a total of 5. I think this is a good deal for me, and I'm willing to accept it.} \\
            }
        }
    }
    & \\ \\

    \theutterance \stepcounter{utterance}  
    & & & \multicolumn{4}{p{0.6\linewidth}}{
        \cellcolor[rgb]{0.9,0.9,0.9}{
            \makecell[{{p{\linewidth}}}]{
                \texttt{\tiny{[GM$\langle$P2]}}
                \texttt{I'm glad we were able to come to a mutually beneficial agreement. I think this trade is a good deal for me as well, since I get the potato, which I value at 3, and one berry, which I value at 3, for a total of 6. However, I'm not sure if this is the optimal deal for both of us, since we could have potentially made a better trade.} \\
\\ 
\texttt{Before we finalize the deal, I'd like to make one more check to see if there's a better option available. Can you confirm that this is the final deal, and that we won't be making any further trades?} \\
            }
        }
    }
     \\ \\

    \theutterance \stepcounter{utterance}  
    & & & \multicolumn{2}{p{0.3\linewidth}}{
        \cellcolor[rgb]{0.9,0.9,0.9}{
            \makecell[{{p{\linewidth}}}]{
                \texttt{\tiny{[GM$|$GM]}}
                \texttt{continue} \\
            }
        }
    }
    & & \\ \\

    \theutterance \stepcounter{utterance}  
    & & & \multicolumn{2}{p{0.3\linewidth}}{
        \cellcolor[rgb]{0.9,0.9,0.9}{
            \makecell[{{p{\linewidth}}}]{
                \texttt{\tiny{[GM$|$GM]}}
                \texttt{I'm glad we were able to come to a mutually beneficial agreement. I think this trade is a good deal for me as well, since I get the potato, which I value at 3, and one berry, which I value at 3, for a total of 6. However, I'm not sure if this is the optimal deal for both of us, since we could have potentially made a better trade.} \\
\\ 
\texttt{Before we finalize the deal, I'd like to make one more check to see if there's a better option available. Can you confirm that this is the final deal, and that we won't be making any further trades?} \\
            }
        }
    }
    & & \\ \\

    \theutterance \stepcounter{utterance}  
    & & \multicolumn{4}{p{0.6\linewidth}}{
        \cellcolor[rgb]{0.9,0.9,0.9}{
            \makecell[{{p{\linewidth}}}]{
                \texttt{\tiny{[P1$\langle$GM]}}
                \texttt{I'm glad we were able to come to a mutually beneficial agreement. I think this trade is a good deal for me as well, since I get the potato, which I value at 3, and one berry, which I value at 3, for a total of 6. However, I'm not sure if this is the optimal deal for both of us, since we could have potentially made a better trade.} \\
\\ 
\texttt{Before we finalize the deal, I'd like to make one more check to see if there's a better option available. Can you confirm that this is the final deal, and that we won't be making any further trades?} \\
            }
        }
    }
    & \\ \\

    \theutterance \stepcounter{utterance}  
    & \multicolumn{4}{p{0.6\linewidth}}{
        \cellcolor[rgb]{0.9,0.9,0.9}{
            \makecell[{{p{\linewidth}}}]{
                \texttt{\tiny{[P1$\rangle$GM]}}
                \texttt{I think we've reached a good agreement, and I'm willing to finalize the deal. However, I'd like to make a secret proposal to confirm that this is the best option for both of us. } \\
\\ 
\texttt{{[}Proposal: 1 watch, 1 berry, 1 potato{]}} \\
            }
        }
    }
    & & \\ \\

    \theutterance \stepcounter{utterance}  
    & & & \multicolumn{2}{p{0.3\linewidth}}{
        \cellcolor[rgb]{0.9,0.9,0.9}{
            \makecell[{{p{\linewidth}}}]{
                \texttt{\tiny{[GM$|$GM]}}
                \texttt{proposal} \\
            }
        }
    }
    & & \\ \\

    \theutterance \stepcounter{utterance}  
    & & & \multicolumn{2}{p{0.3\linewidth}}{
        \cellcolor[rgb]{0.9,0.9,0.9}{
            \makecell[{{p{\linewidth}}}]{
                \texttt{\tiny{[GM$|$GM]}}
                [0, 1, 1, 1]
            }
        }
    }
    & & \\ \\

    \theutterance \stepcounter{utterance}  
    & & \multicolumn{4}{p{0.6\linewidth}}{
        \cellcolor[rgb]{0.9,0.9,0.9}{
            \makecell[{{p{\linewidth}}}]{
                \texttt{\tiny{[GM$\rangle$P2]}}
                \texttt{The other player has submitted a secret proposal. It is now your turn to submit a proposal of your own.} \\
            }
        }
    }
    & \\ \\

    \theutterance \stepcounter{utterance}  
    & & & \multicolumn{4}{p{0.6\linewidth}}{
        \cellcolor[rgb]{0.9,0.9,0.9}{
            \makecell[{{p{\linewidth}}}]{
                \texttt{\tiny{[GM$\langle$P2]}}
                \texttt{{[}Proposal: 1 berry, 1 potato, 1 plier{]}} \\
            }
        }
    }
     \\ \\

    \theutterance \stepcounter{utterance}  
    & & & \multicolumn{2}{p{0.3\linewidth}}{
        \cellcolor[rgb]{0.9,0.9,0.9}{
            \makecell[{{p{\linewidth}}}]{
                \texttt{\tiny{[GM$|$GM]}}
                \texttt{proposal} \\
            }
        }
    }
    & & \\ \\

    \theutterance \stepcounter{utterance}  
    & & & \multicolumn{2}{p{0.3\linewidth}}{
        \cellcolor[rgb]{0.9,0.9,0.9}{
            \makecell[{{p{\linewidth}}}]{
                \texttt{\tiny{[GM$|$GM]}}
                [1, 1, 0, 1]
            }
        }
    }
    & & \\ \\

    \theutterance \stepcounter{utterance}  
    & & & \multicolumn{2}{p{0.3\linewidth}}{
        \cellcolor[rgb]{0.9,0.9,0.9}{
            \makecell[{{p{\linewidth}}}]{
                \texttt{\tiny{[GM$|$GM]}}
                [[0, 1, 1, 1], [1, 1, 0, 1]]
            }
        }
    }
    & & \\ \\

    \theutterance \stepcounter{utterance}  
    & \multicolumn{4}{p{0.6\linewidth}}{
        \cellcolor[rgb]{0.95,0.95,0.95}{
            \makecell[{{p{\linewidth}}}]{
                \texttt{\tiny{[P2$\rightarrow$P1]}}
                \texttt{Post{-}episode Reflexion (not sent to GM):} \\
\texttt{New plan: } \\
\\ 
\texttt{In my previous attempts, I failed to achieve a mutually beneficial trade that would maximize my score. I got stuck in a cycle of negotiations, trying to get the berries and the plier, but ultimately failing to secure a good deal. } \\
\\ 
\texttt{What went wrong was that I didn't take into account the other player's high value for the berries and the plier's low value to them. I should have focused on getting the berries and the watch, which would have given me a higher score. } \\
\\ 
\texttt{However, upon further reflection, I realize that I made another critical mistake. I didn't make a secret proposal until the very end, after the other player had already agreed to the trade. This meant that I didn't have a clear plan for how to formalize the trade and ensure that both of us would get the items we wanted.} \\
\\ 
\texttt{Additionally, I also failed to follow the format for making a secret proposal, which is to use square brackets and the correct format. } \\
\\ 
\texttt{Next time, I will take a more strategic approach by considering the other player's values and trying to find a mutually beneficial trade that maximizes my score. I will also pay closer attention to the other player's counteroffers and be more willing to compromise. } \\
\\ 
\texttt{Here's a revised plan:} \\
\\ 
\texttt{1. Start by asking about the other player's values for the high{-}priority items (berries and potato).} \\
\texttt{2. Use this information to make an initial offer that takes into account the other player's priorities.} \\
\texttt{3. Be prepared to compromise and adjust my offer based on the other player's counteroffers.} \\
\texttt{4. Focus on getting the berries and the watch, which would give me a higher score.} \\
\texttt{5. Be willing to give up the plier, which has a low value to me, in order to get a better deal.} \\
\texttt{6. Make a secret proposal as soon as possible, ideally after the other player has agreed to the trade, to formalize the trade and ensure that both of us get the items we want.} \\
\texttt{7. Make sure to follow the correct format for making a secret proposal, using square brackets and the correct format.} \\
\texttt{8. Before making a secret proposal, I will make sure to confirm with the other player that we have a mutually beneficial trade, and that both of us are satisfied with the terms of the trade.} \\
\\ 
\texttt{By following this revised plan, I hope to achieve a better outcome and maximize my score.} \\
            }
        }
    }
    & & \\ \\

    \theutterance \stepcounter{utterance}  
    & & & \multicolumn{4}{p{0.6\linewidth}}{
        \cellcolor[rgb]{0.95,0.95,0.95}{
            \makecell[{{p{\linewidth}}}]{
                \texttt{\tiny{[P1$\rightarrow$P2]}}
                \texttt{Post{-}episode Reflexion (not sent to GM):} \\
\texttt{New plan: } \\
\\ 
\texttt{In my last attempt, I failed to achieve a mutually beneficial trade that would have maximized my score. I got stuck in a cycle of negotiations, trying to get the berries and the potato, but ultimately failing to secure a deal. } \\
\\ 
\texttt{Upon reflection, I realize that I should have focused more on the value function of the other player. I knew that the other player valued the berries highly, but I didn't use this information effectively to my advantage. } \\
\\ 
\texttt{Next time, I will try to understand the other player's priorities and value function more clearly, and use this information to make more informed decisions. I will also try to be more flexible and open to different trade{-}offs, rather than getting stuck on a single goal. } \\
\\ 
\texttt{Specifically, I will pay closer attention to the other player's responses and try to understand what they are willing to give up and what they are holding out for. I will also try to make more creative and mutually beneficial proposals, rather than simply trying to get what I want. } \\
\\ 
\texttt{To achieve this, I will start by making a more detailed analysis of the other player's value function, and try to identify their priorities and potential trade{-}offs. I will also try to make more open{-}ended and exploratory proposals, rather than simply trying to get a specific item or set of items. } \\
\\ 
\texttt{For example, I could have proposed a trade that would have given the other player the berries and the watch, while also giving me the potato and the plier. This would have been a more balanced and mutually beneficial deal, and would have given me a higher score. } \\
\\ 
\texttt{Next time, I will try to think more strategically and make more informed decisions based on the other player's priorities and value function.} \\
\\ 
\texttt{Additionally, I will pay closer attention to the format of the secret proposal, and make sure to use the correct format when making a secret proposal. I will also avoid using square brackets when communicating with the other player, as this will be interpreted as a secret proposal. } \\
\\ 
\texttt{For example, instead of saying "{[}Proposal: 1 berry, 1 potato{]}", I will say "I propose that we trade 1 berry for 1 potato". This will ensure that my message is clear and follows the correct format.} \\
\\ 
\texttt{Furthermore, I will try to make a secret proposal that is more balanced and takes into account the other player's priorities and value function. I will also try to make a proposal that is more creative and mutually beneficial, rather than simply trying to get a specific item or set of items.} \\
\\ 
\texttt{For example, I could propose a trade that would give the other player the berries and the plier, while also giving me the potato and the watch. This would be a more balanced and mutually beneficial deal, and would give me a higher score.} \\
\\ 
\texttt{Next time, I will try to think more strategically and make more informed decisions based on the other player's priorities and value function. I will also try to be more flexible and open to different trade{-}offs, rather} \\
            }
        }
    }
     \\ \\

\end{supertabular}
}

\end{document}
